% concatenated from 0417_Normal_approximation_for_iterated_inner_functions.tex
\documentclass[11pt, reqno]{amsart}
%\documentclass[reqno]{amsart}
%\documentclass[11pt]{article}
%\documentclass[12pt]{amsart}

%--------------------------------------Packages---------------------------------
\usepackage[colorlinks]{hyperref}
\usepackage{amssymb}
\usepackage{mathrsfs}
\usepackage[all,cmtip]{xy}
\usepackage{cite}
\usepackage{color}
\usepackage{tikz-cd}
\usepackage{framed}
\usepackage{enumitem}

%----------------------------------page setup------------------------------
\textheight=22.5cm
\textwidth=15.0cm
\oddsidemargin=0.4in
\evensidemargin=0.4in
\topmargin=-.5cm
\footskip=1.5cm
\renewcommand{\baselinestretch}{1.3}

%--------------------------------Theorem Environment----------------------------
\newtheorem{theorem}{Theorem}
\newtheorem{corollary}[theorem]{Corollary}
\newtheorem{lemma}[theorem]{Lemma}
\newtheorem{proposition}[theorem]{Proposition}
\newtheorem{example}[theorem]{Example}
\newtheorem{remark}[theorem]{Remark}
\newtheorem{question}[theorem]{Question}
\theoremstyle{definition}
\newtheorem{definition}[theorem]{Definition}
\counterwithout{equation}{section}
\newtheorem*{Brown-Eagleson Theorem}{The Brown-Eagleson Theorem}



%----------------------------------newcommand--------------------
\newcommand{\D}{\mathbb D}
\newcommand{\R}{\mathbb R}
\newcommand{\C}{\mathbb C}
\newcommand{\F}{\mathscr F}
\newcommand{\T}{\mathbb T}
\newcommand{\E}{\mathbb E}
\newcommand{\PP}{\mathbb P}
\newcommand{\cal}{\mathcal}
\newcommand{\ol}{\overline}
\newcommand{\Ran}{{\rm Ran}\,}
\newcommand{\Ker}{{\rm Ker}\,}
\newcommand{\la}{\langle}
\newcommand{\ra}{\rangle}
\newcommand{\HH}{\mathbb H}
\newcommand{\one}{{\bf 1}}
\newcommand{\M}{\mathcal M}
\newcommand{\X}{\mathfrak{X}}
\newcommand{\XX}{\mathcal{X}}
\newcommand{\sgn}{{\rm sgn}\,}


\allowdisplaybreaks



%------------------------------------article information-----------------------
\title{Normal approximation for iterated inner functions}

\author{Yukun Chen}
\address{Yukun Chen: School of Mathematics and Statistics, Wuhan University, Wuhan 430072, China}
\email{yukunchen@whu.edu.cn}


\author{Xiangdi Fu}
\address{Xiangdi Fu: School of Fundamental Physics and Mathematical Sciences, HIAS, University of Chinese Academy of Sciences, Hangzhou, 310024, China}
\email{xdfu@ucas.ac.cn}


\author{Zhaofeng Lin}
\address{Zhaofeng Lin: School of Fundamental Physics and Mathematical Sciences, HIAS, University of Chinese Academy of Sciences, Hangzhou, 310024, China}
\email{linzhaofeng@ucas.ac.cn}

\author{Yanqi Qiu}
\address{Yanqi Qiu: School of Fundamental Physics and Mathematical Sciences, HIAS, University of Chinese Academy of Sciences, Hangzhou, 310024, China}
\email{yanqiqiu@ucas.ac.cn}





\keywords{Berry--Ess\'{e}en theorem, martingale central limit theorem, inner functions}
%\thanks{}



\makeatletter
\@namedef{subjclassname@2020}{\textup{2020} Mathematics Subject Classification}
\makeatother
\subjclass[2020]{Primary: 30J05; Secondary: 60F05, 37F10.}


\begin{document}
	%--------------------------------------Abstract----------------------------------------------
	\begin{abstract}
	 A Berry--Ess\'{e}en theorem for linear combinations of iterates of an inner function is obtained. Our proof, which is based an elementary transfer argument and classical results in martingale theory, also leads to a simple proof of Nicolau and Soler i Gibert's central limit theorem for inner functions.
	\end{abstract}
	\maketitle
	
	%--------------------------------------Contents---------------------------------------------- 

	

	\section{Introduction}
	Let $\C$ be the complex plane, and let $\D= \{z\in \C: |z|<1\}$ be the open unit disc. Let $\T=\{z\in \C: |z|=1\}$ be the unit circle, equipped with the normalized Lebesgue measure $m$. An analytic function $f:\D \to \D$ is called {\it inner} if its radial limits $$f(\omega):= \lim_{r\to 1} f(r\omega)$$ have modulus one at almost every $\omega\in \T$. Therefore, by taking radial limits, any inner function $f$ induces a boundary map from $\T$ into itself defined at almost every point. This boundary map will also be denoted by $f$. 
	
	The study of inner functions is a fundamental topic in complex analysis, and a comprehensive account of their properties and related theories can be found in classical monographs; see \cite{Gar81} for example. It is well known that if $f(0)=0$, then the corresponding boundary map preserves the normalized Lebesgue measure $m$, meaning $m(f^{-1}(E)) = m (E)$ for all measurable set $E\subset \T$. This measure-preserving property has motivated extensive research into their dynamical behavior \cite{DM91,Pom81}, which has in turn drawn considerable attention to the iteration of inner functions \cite{Nic22,NS22,NS24}.

	Throughout this note, we denote by $f^{\circ n}$ the $n$-th iterate of $f$, and by $\mathcal{N}(0,\sigma^2)$ the real normal distribution with mean zero and variance $\sigma^2$. We also denote by $\mathcal{CN}(0,\sigma^2)$ the circularly symmetric complex normal distribution with mean zero and variance $\sigma^2$ (equivalently, its real and imaginary parts are independent $\mathcal{N}(0,\sigma^2/2)$ variables).

	The main result of this note is the following Berry--Ess\'{e}en type theorem for iterated inner functions. 


\begin{theorem}\label{thm-BE-inner}
	Let $f$ be an inner function with $f(0)=0$ which is not a rotation. Let $(a_n)_{n=1}^\infty$ be a sequence of complex numbers satisfying
	\begin{align}\label{eq-Lindeberg-cond}\sum_{n=1}^\infty |a_n|^2 = \infty,\quad \text{and} \quad \lim_{N\to \infty} \frac{|a_N|^2}{\sum_{n=1}^N |a_n|^2} = 0,
	\end{align} 
	and set $$Z_N:= \frac{1}{\sigma_N}\sum_{n=1}^N a_n f^{\circ n},\quad \text{and }\,\, \sigma_N:= \bigg\| \sum_{n=1}^N a_n f^{\circ n }\bigg\|_{L^2(\T)}.$$ Then for any $\delta\geq 1$, there exist positive constants $C(\delta,f)$ (depending only on $\delta$ and $f$), such that the following Cram\'er-Wold error estimate holds: 
		\begin{equation}
			\begin{split}
		&\sup_{x\in \R,\, \alpha \in \T} \bigg| \PP\Big( {\rm Re}\, \big( \alpha Z_N\big) \leq x \Big) - \PP\Big( \cal{N}(0,1/2) \leq x \Big)\bigg|  \\ 
		& \,\leq \,C(\delta,f) \cdot \bigg(\sum_{n=1}^N |a_n|^{4} \bigg)^{\frac{1+\delta}{6+4\delta}}  \bigg(\sum_{n=1}^N |a_n|^{2} \bigg)^{-\frac{1+\delta}{3+2\delta}} + 2.7 \cdot \frac{|\lambda|}{1-|\lambda|} \frac{\big|\sum_{n=1}^N \lambda^{N-n} a_n\big|}{\sqrt{\sum_{n=1}^N |a_n|^2}}.
			\end{split}\label{eq-rate-an}
		\end{equation}
	\end{theorem}

	We point out that Theorem \ref{thm-BE-inner} provides a refinement of the Central Limit Theorem (CLT) for inner functions (Theorem \ref{thm-NS} below) recently established by A. Nicolau and O. Soler i Gibert \cite{NS22,NS24}. Indeed, if the sequence \((a_n)_{n=1}^\infty \) satisfies the conditions in \eqref{eq-Lindeberg-cond}, then the quantity in \eqref{eq-rate-an} tends to zero for any fixed \(\delta \geq 1\), thereby establishing the convergence in distribution: $$\frac{1}{\sigma_N} \sum_{n=1}^N a_n f^{\circ n} \stackrel{d}{\longrightarrow} \cal{CN}(0,1).$$ Moreover, in the constant coefficient case $a_n \equiv 1$, the bound \eqref{eq-rate-an} provides a decay rate of $O(N^{-\frac{1}{4} + \varepsilon})$ for any $\varepsilon > 0$. An interesting problem is to determine the optimal rate of decay of the errors arising in the convergence in distribution in the central limit theorem for inner functions. 
	
	%In fact, we also observed that if $a_n \equiv 1$ and $f''(0)=0$, then the quantity in \eqref{eq-CW-error} decays at the rate $O(n^{-1/2}\log n)$.


	The proof of Theorem \ref{thm-BE-inner} is based on an elementary transfer argument together with classical results in martingale theory. In particular, our method also yields a short proof of Nicolau and Soler i Gibert's CLT. 

	\noindent\textbf{The structure of this paper:} Since the martingale CLT is considerably more elementary and more familiar to a general audience than its Berry--Ess\'{e}en counterpart, we prefer to first present our short proof of Nicolau and Soler i Gibert's CLT, and then demonstrate how a simple modification of the argument leads to Theorem \ref{thm-BE-inner}.


	\section{A simple proof of the central limit theorem for inner functions}


	In this section, we present a simple proof of the following

		\begin{theorem}[Nicolau and Soler i Gibert]\label{thm-NS}
		Let $f$ be an inner function with $f(0)=0$ which is not a rotation. Let $(a_n)_{n=1}^\infty$ be a sequence of complex numbers satisfying the conditions in \eqref{eq-Lindeberg-cond}, and set $$\sigma_N:= \bigg\| \sum_{n=1}^N a_n f^{\circ n }\bigg\|_{L^2(\T)}.$$ Then the following convergence in distribution holds: 
		\begin{align}\label{eq-convergence-distribution}\frac{1}{\sigma_N} \sum_{n=1}^N a_n f^{\circ n }\stackrel{d}{\longrightarrow} \cal{CN}(0,1).
		\end{align}
	\end{theorem}

	We note that the normal approximation \eqref{eq-convergence-distribution} was first established in \cite{NS22} under certain conditions on the size of the coefficient sequence $(a_n)_{n=1}^\infty$, and subsequently refined in \cite{NS24} to the sharp form stated in Theorem \ref{thm-NS}. The original proof relies on the theory of Aleksandrov-Clark measures as the main technical tool to establish the pointwise convergence of the characteristic functions.



	\subsection{The reverse martingale differences}
	
	Let $f$ be an inner function with $f(0)=0$, which is not a rotation. We denote by $\F_0$ the $\sigma$-algebra consisting of all Lebesgue measurable sets on $\mathbb{T}$, and let $\F_n$ be the complete sub-$\sigma$-algebra of $\F_0$ generated by $f^{\circ n}$. Clearly, the sequence $(\F_n)_{n=0}^\infty$ is decreasing: $$\F_0 \supset  \F_1 \supset \F_2 \supset \cdots.$$ To simplify notation, we set $\lambda := \overline{f'(0)}$ and $\mu:=\ol{f''(0)}/2$. Note that (by the Schwarz--Pick lemma)  $$|\lambda|<1, \quad \text{and } \, |\mu|\leq 1-|\lambda|^2.$$
	
	Now we define 
	\begin{align}\label{eq-def-Yn}
		Y_n:= f^{\circ n}-\lambda f^{\circ (n+1)}.
	\end{align} It turns out that $(Y_n)_{n=1}^\infty$ naturally forms a \emph{reverse} martingale difference sequence. The details are presented in the following lemma.


	\begin{lemma}\label{lem-Yn}
		The sequence $(Y_n)_{n=1}^\infty$ forms a reverse martingale difference sequence with respect to $(\F_n)_{n=0}^\infty$, that is, $$Y_n \text{ is $\F_n$-measurable}, \,\, \text{and }\,\, \E[Y_n|\F_{n+1}] = 0.$$
		Moreover, we have the following identities:
		\begin{enumerate}[label=(\roman*),font=\normalfont]
			\item \label{item:|Yn|2} $\E\big[|Y_n|^2 | \F_{n+1}\big]=\E\big[|Y_n|^2\big] = 1-|\lambda|^2.$
			\item \label{item:Yn2} $\E\big[Y_n^2| \F_{n+1}\big]=\mu f^{\circ (n+1)}$.
			\item \label{item:reYn2} For any $\alpha \in \mathbb C$, $$\E\big[({\rm Re}\, \alpha Y_n)^2 |\mathscr F_{n+1}\big]=\frac{1-|\lambda|^2}{2}|\alpha|^2 + \frac{{\rm Re}\,(\alpha^2 \mu f^{\circ (n+1)})}{2}.$$
		\end{enumerate}
		
	\end{lemma}
	
		The proof of Lemma \ref{lem-Yn} requires only the following standard result: if $\eta$ is an inner function vanishing at the origin, then the collection $$\{\eta^k : k\in \mathbb Z \}$$ constitutes an orthonormal basis for $L^2(\F(\eta)):=\{f\in L^2(\T): \text{$f$ is $\F(\eta)$-measurable}\}$, where $\F(\eta)$ is the complete sub $\sigma$-algebra generated by $\eta$. Since the conditional expectation $\E\big[ \cdot | \F(\eta)\big]$ acts as the orthogonal projection from  $L^2(\T)$ onto $L^2(\F(\eta))$, it holds that $$\E\big[ h| \F(\eta)\big] = \sum_{k\in \mathbb Z} \, \big\la h, \eta^k \big\ra \cdot \eta^k,\quad h\in L^2.$$ The rest of the proof of Lemma \ref{lem-Yn} consists of standard calculations. To keep the exposition focused, we have postponed it (and the proof of Lemma \ref{lem-transfer} below) to the end of this note.

		\subsection{The transfer argument}
		
		The second key ingredient is the following elementary identity, which rewrites a linear combination of iterates of inner functions as a linear combination of reverse martingale differences plus a remainder term.


	\begin{proposition}\label{lem-inner=Yn+remainder}
	For any complex sequence $(a_n)_{n=1}^\infty$, and any positive integer $N$, it holds that: 
		\begin{align}\label{eq-sum-to-martingale}
			\sum_{n=1}^N a_n f^{\circ n}=\sum_{n=1}^N b_n Y_n + \lambda b_N f^{\circ (N+1)}, 
		\end{align}
		where the coefficients $b_n$ are given by 
		\begin{align}\label{eq-def-bn}
			b_n=\sum_{k=1}^n \lambda^{n-k} a_k, \quad n\geq 1.
		\end{align}
	\end{proposition}

One can easily show that if the original coefficients satisfy the conditions in \eqref{eq-Lindeberg-cond}, then so do the resulting coefficients:
	
	\begin{lemma}\label{lem-transfer}
		Let $a=(a_n)_{n=1}^\infty$ be any complex sequence, and let $b=(b_n)_{n=1}^\infty$ be the associated sequence defined as in \eqref{eq-def-bn}. Then the following claims hold true: 
		\begin{enumerate}[label=(\roman*),font=\normalfont]
			\item \label{item:keep-Lindberg}If $$\lim_{N\to \infty} \frac{\max_{1\leq n \leq N} |a_n|^2}{\sum_{n=1}^N |a_n|^2} = 0,$$ then $$\lim_{N\to \infty} \frac{\max_{1\leq n \leq N} |b_n|^2}{\sum_{n=1}^N |b_n|^2} =0.$$
			\item \label{item:keep-no-ell2} If $a\notin \ell^2$, then $b\notin \ell^2$.
			\item \label{item:keep-ell2}If $a\in \ell^2$, then $b\in \ell^2$ and $$\frac{1}{1+|\lambda|}\cdot \|a\|_{\ell^2} \leq \|b\|_{\ell^2} \leq \frac{1}{1-|\lambda|}\cdot \|a\|_{\ell^2}.$$
		\end{enumerate}
	\end{lemma}

	Note that if $a\notin \ell^2$, then $$\lim_{N\to \infty} \frac{|a_N|^2}{\sum_{n=1}^N |a_n|^2}=0 \text{\, if and only if \,}\lim_{N\to \infty} \frac{\max_{1\leq n\leq N}|a_n|^2}{\sum_{n=1}^N |a_n|^2}=0.$$ Moreover, in the case $a\in \ell^2$, one can establish the following inequalities:
	\begin{align}\label{eq-2-norm-eqvalence}
		\frac{1-|\lambda|}{1+|\lambda|} \cdot \|a\|^2_{\ell^2} \leq \bigg\|\sum_{n=1}^\infty a_n f^{\circ n}\bigg\|^2_2 \leq \frac{1+|\lambda|}{1-|\lambda|}\cdot \|a\|^2_{\ell^2}.
	\end{align}
	Indeed, for any $a\in \ell^2$, the associated sequence $b$ defined by \eqref{eq-def-bn} is also square summable. Letting $N\to \infty$ in \eqref{eq-sum-to-martingale} yields $$\sum_{n=1}^\infty a_n f^{\circ n} = \sum_{n=1}^\infty b_n Y_n.$$ Since the reverse martingale differences $Y_n$ are pairwise orthogonal and $\E\big[ |Y_n|^2\big] = 1-|\lambda|^2$, we conclude that $$\bigg\|\sum_{n=1}^\infty a_n f^{\circ n}\bigg\|^2_2 = \big( 1- |\lambda|^2\big) \cdot \|b\|^2_2.$$ Combining this with Lemma \ref{lem-transfer}-\ref{item:keep-ell2}, we obtain the inequalities in \eqref{eq-2-norm-eqvalence}. Moreover, since $Y_n \perp f^{\circ (N+1)}$ for each $1\leq n \leq N$, it follows from \eqref{eq-sum-to-martingale} that 
	\begin{align}\label{eq-rhoN-sigmaN} 
		\bigg\|\sum_{n=1}^N a_n f^{\circ n}\bigg\|^2_2 = (1-|\lambda|^2)\sum_{n=1}^N |b_n|^2  + |\lambda|^2 |b_N|^2.
	\end{align}
	
	It should be pointed out that \eqref{eq-2-norm-eqvalence} was established in \cite[Theorem 9]{NS22} through properties of Toeplitz matrices. 
	
	\subsection{A weak law of large numbers for inner functions}

	\begin{lemma}\label{lem-weak-law}
		Let $f$ be an inner function with $f(0)=0$ which is not a rotation. Suppose the complex sequence $(a_n)_{n=1}^\infty$ satisfies $$ S_N:=\sum_{n=1}^N |a_n| \to \infty, \quad \text{and}\quad  \frac{|a_N|}{S_N}\to 0.$$ Then $\frac{1}{S_N} \sum_{n=1}^N a_n f^{\circ n} \stackrel{p}{\longrightarrow} 0$, i.e.  converges to zero in probability. 
	\end{lemma}
	\begin{proof} 
		We only need to prove $\frac{1}{S_N} \sum_{n=1}^N a_n f^{\circ n}$ converges to zero in $L^2$. In view of \eqref{eq-2-norm-eqvalence}, it suffices to prove $$\lim_{N\to \infty} \frac{\sum_{n=1}^N |a_n|^2}{S^2_N}=0.$$ Without loss of generality, we may assume $a_n\neq 0$ for each $n\geq 1$ and $\sum_{n=1}^\infty |a_n|^2 = \infty$. Then it follows from the Stolz--Ces\`aro theorem that $$\limsup_{N\to \infty} \frac{\sum_{n=1}^N |a_n|^2}{S^2_N} \leq  \limsup_{N\to \infty} \frac{|a_N|^2}{S_N^2-S^2_{N-1}}= \limsup_{N\to \infty}\frac{|a_N|}{ S_N+S_{N-1}}=0.$$
		The proof is completed. 
	\end{proof}

	\subsection{A central limit theorem for martingale triangular arrays}

	 We need to employ a classical result due to B. M. Brown and G. K. Eagleson \cite[Corollary 1]{BE71}. For the reader's convenience, we provide its precise statement here.

	\begin{Brown-Eagleson Theorem}\label{thm-BE71}
		Consider the following triangular array of real-valued random variables on a probability space $(\Omega, \cal F, \mathbb P)$:
		\begin{align}\label{eq-martingale-array}
		\begin{pmatrix}
 			X_{1,1} & 0 & 0 & \cdots \\
			X_{2,1} & X_{2,2} & 0 & \cdots \\
			X_{3,1} & X_{3,2} & X_{3,3} & \cdots \\
			\vdots & \vdots & \vdots & \ddots
		\end{pmatrix}
	\end{align}
		Let $\{\mathcal{F}_{N,n}: 0\leq n \leq N\}$ be a family of $\sigma$-algebras such that for each $N\geq 1$ fixed, 
		\begin{align}\label{eq-condition-incresing-chain}
			\cal F_{N,0} \subset \cal F_{N,1} \subset \cdots \subset \cal F_{N,N} \subset \cal F, 
		\end{align} 
		and for any $1\leq n\leq N$, 
		\begin{align}\label{eq-condition-martingale} 
			X_{N,n} \text{ is $\cal F_{N,n}$-measurable}, \quad \E\big[ X_{N,n} | \cal F_{N,n-1}\big]=0.
		\end{align}
		Define
		$$
			V^2_N:= \sum_{n=1}^N \E\big[X_{N,n}^2| \cal F_{N,n-1}\big],\quad  m_N:= \max_{1\leq n\leq N} \E\big[X_{N,n}^2| \cal F_{N,n-1}\big].
		$$ 
		If the following three conditions are satisfied as $N\to \infty$: 
		\begin{itemize}
			\item [(C1)] $m_N \stackrel{p}{\longrightarrow} 0$;
			\item [(C2)] $V^2_N \stackrel{p}{\longrightarrow} \sigma^2$ (constant);
			\item [(C3)] For any $\varepsilon>0$, $$\sum_{n=1}^N \E[X_{N,n}^2 I\big(|X_{N,n}| \geq \varepsilon\big)| \cal F_{N,n-1}] \stackrel{p}{\longrightarrow}  0,$$
		\end{itemize}
		 then the $N$-th row sum of \eqref{eq-martingale-array} converges in law to $\cal N(0,\sigma^2)$ as $N\to \infty$: $$X_{N,1}+ \cdots + X_{N,n}\stackrel{d}{\longrightarrow} \cal N(0,\sigma^2).$$
	\end{Brown-Eagleson Theorem}

	For further generalizations of this result, we refer the reader to \cite{Eag75}.
	


	\subsection{Concluding the proof}

	\begin{proof}[Proof of Theorem \ref{thm-NS}]
		Let $Y_n$ and $b_n$ be defined as in \eqref{eq-def-Yn} and \eqref{eq-def-bn}, respectively. Then by Lemma \ref{lem-transfer}, 
		\begin{align}\label{eq-bn-Lindberg}
			\sum_{n=1}^\infty |b_n|^2=\infty,\quad {\rm and} \quad \lim_{N\to \infty} \frac{|b_N|^2}{\sum_{n=1}^N |b_n|^2}= 0.
		\end{align}
		Combining this with Lemma \ref{lem-Yn}-\ref{item:|Yn|2} and \eqref{eq-sum-to-martingale}, one has $$\lim_{N\to \infty} \frac{\sigma^2_N}{\rho_N^2} = 1-|\lambda|^2, \quad {\rm where}\,\, \rho_N:= \sqrt{\sum_{n=1}^N |b_n|^2}.$$ In virtue of the identity \eqref{eq-sum-to-martingale}, to prove $$\frac{1}{\sigma_N}\sum_{n=1}^N a_nf^{\circ n}\stackrel{d}{\longrightarrow}\cal{CN}(0,1),$$ it suffices to show
		$$\frac{1}{\rho_N}\sum_{n=1}^N b_n Y_n \stackrel{d}{\longrightarrow} \mathcal{CN}(0,1-|\lambda|^2).$$ Furthermore, by circular symmetry, we only need to establish that $${\rm Re}\,\bigg(\frac{\alpha}{\rho_N}\sum_{n=1}^N b_n Y_n\bigg) \stackrel{d}{\longrightarrow} \mathcal N\bigg(0,\frac{1-|\lambda|^2}{2}\bigg),$$ for any fixed $\alpha\in \T$.
		
		
		Now we set $$X_{N,n}:= {\rm Re}\, \bigg( \frac{\alpha \,b_{N+1-n} Y_{N+1-n}}{\rho_N}\bigg),\quad 1\leq n\leq N,$$ and $$\mathcal{F}_{N,n} := \mathscr{F}_{N-n+1},\quad 0\leq n \leq N.$$ The proof proceeds by applying the Brown--Eagleson theorem to the triangular array \eqref{eq-martingale-array}. Conditions \eqref{eq-condition-incresing-chain} and \eqref{eq-condition-martingale} are clearly met, so it suffices to check conditions (C1)-(C3) there.
	

		By Lemma \ref{lem-Yn}-\ref{item:reYn2} and \eqref{eq-bn-Lindberg},
		\begin{align*}
			m_N&=\frac{1}{\rho^2_N} \max_{1\leq n \leq N} \bigg(\frac{1-|\lambda|^2}{2} |b_n|^2+ \frac{{\rm Re}\,(\alpha^2 \mu \,b_n^2f^{\circ (n+1)})}{2}\bigg)\\
			&\leq \frac{\max_{1\leq n\leq N} |b_n|^2}{\rho^2_N}\to 0, \quad \text{as $N\to \infty$}.
		\end{align*}
		This implies (C1). Moreover, we have 
		\begin{align} 
			V_N^2&= \frac{1}{\rho^2_N}\sum_{n=1}^N \bigg(\frac{1-|\lambda|^2}{2} |b_n|^2+ \frac{{\rm Re}\,(\alpha^2 \mu \,b_n^2f^{\circ (n+1)})}{2}\bigg) \nonumber \\
			&=\frac{1-|\lambda|^2}{2}+ \frac{1}{\rho^2_N}\sum_{n=1}^N \bigg(\frac{{\rm Re}\,(\alpha^2 \mu \,b_n^2f^{\circ (n+1)})}{2}\bigg). \label{eq-VN-value}
		\end{align}
		An application of Lemma \ref{lem-weak-law} implies $$\frac{1}{\rho^2_N} \sum_{n=1}^N b_n^2 f^{\circ n} \stackrel{p}{\longrightarrow} 0,$$ and hence the second term of $V_N^2$ converges to zero in probability. This shows (C2) holds true with $\sigma^2 = \frac{1-|\lambda|^2}{2}$. It remains to prove (C3). Note that $$\max_{1\leq n\leq N} \frac{{\rm Re}\, (\alpha b_n Y_n)}{\rho_N} \leq  \big(1+|\lambda|\big)\frac{\max_{1\leq n\leq N}|b_n|}{\rho_N}\to 0, \quad N\to \infty.$$ So for any fixed $\varepsilon>0$, we have
		\begin{align*}
			&\sum_{n=1}^N \E[X_{N,n}^2 I(|X_{N,n}|>\varepsilon)| \cal F_{N,n-1}]\\
			=&\frac{1}{\rho^2_N} \sum_{n=1}^N \E\bigg[\bigg({\rm Re}\, (\alpha b_n Y_n)\bigg)^2 I\bigg(\frac{{\rm Re}\, (\alpha b_n Y_n)}{\rho_N}> \varepsilon \bigg) \bigg| \F_{n+1}\bigg]\\
			\leq& \frac{1}{\rho^2_N} \sum_{n=1}^N \E\bigg[\bigg({\rm Re}\, (\alpha b_n Y_n)\bigg)^2 I\bigg(\max_{1\leq n\leq N}\frac{{\rm Re}\, (\alpha b_n Y_n)}{\rho_N}> \varepsilon\bigg) \bigg| \F_{n+1}\bigg]\\
			=&0
		\end{align*}
		provided that $N$ is sufficiently large. This proves (C3) and completes the proof of Theorem \ref{thm-NS}.
\end{proof}

	\noindent \textbf{Remark.} Parallel to Theorem \ref{thm-NS}, Nicolau and Soler i Gibert also established a CLT for tails: if $a\in \ell^2$ and $$\lim_{N\to \infty} \frac{|a_N|^2}{\sum_{n=N}^\infty |a_n|^2}=0,$$ then with $\sigma(N):= \big\|\sum_{n=N}^\infty a_n f^{\circ n}\big\|_2$, it holds that $$\frac{1}{\sigma(N)}\sum_{n=N}^\infty a_n f^{\circ n}\stackrel{d}{\longrightarrow} \mathcal{CN}(0,1),\quad \text{ as $N\to \infty$}.$$ The arguments used here can also be adapted to prove this result. In this case, the proof is even simpler: it follows directly from the central limit theorem for reverse martingales \cite{Loy69}, without requiring the triangular array version. 




	



	





	\section{Proof of Theorem \ref{thm-BE-inner}}




We apply the same strategy as in the proof of Theorem \ref{thm-NS}, substituting the Brown--Eagleson Theorem with a deeper result of E. Haeusler \cite[Theorem 1]{Hae88}. The latter asserts that: for any $\delta>0$, there exists a constant $C$ depending only on $\delta$, such that
\begin{align*}
	&\sup_{x\in \R}\bigg| \PP\Big( \sum_{n=1}^N X_n\leq x \Big) - \PP \Big(\cal N(0,1) \leq x \Big)\bigg| \\
	\leq & C\cdot  \bigg( \sum_{n=1}^N \E\big[ |X_n|^{2+2\delta}\big]+ \E\bigg[\, \bigg|\sum_{n=1}^N \E\big[ X_n^2 | \cal F_{n-1}\big] -1\bigg|^{1+\delta}\bigg]\bigg)^{\frac{1}{3+2\delta}}
\end{align*}
holds for any real-valued martingale difference sequence $(X_n)_{n=1}^N$ of length $N$ (with respect to the filtration $(\cal F_n)^N_{n=0}$).

We also need the following lemma. 
\begin{lemma}\label{lem-normal-ball-probability}
	Let $\Phi(x)$ be the distribution function of $\mathcal N(0,1/2)$. Then
	$$
	\sup_{x\in \R}\bigg| \Phi\big(px+q\big) - \Phi\big(x\big)\bigg|
	\leq 1.35 \cdot (|p-1|+|q|),
	$$
	for any $1/2\leq p\leq 3/2$ and $-1\leq q \leq 1$.
\end{lemma}

\begin{proof} 
	By the mean value theorem, for any $x\in \R$, there exists $\xi$ between $x$ and $px+q$, such that
	$$
	\Big|\Phi\big(px+q\big)- \Phi(x)\Big|
	= \frac{|(p-1)x+q|}{\sqrt{\pi}} \,e^{-\xi^2}.
	$$
	Observe that
	$$
	|\xi|\geq \inf \big\{|px+q|: 1/2 \leq p \leq 3/2,\ -1\leq q \leq 1\big\}
	=\max\big\{0,|x|/2 -1 \big\}=: m(x),
	$$
	and hence
	$$
	C:= \frac{1}{\sqrt{\pi}}\max \bigg\{ \sup_{x\in \R} |x|\cdot e^{-m(x)^2}, \, 1\bigg\}\approx 1.348... <1.35 
	$$
	provides a desired upper bound.
\end{proof}

\begin{proof}[Proof of Theorem \ref{thm-BE-inner}]
We use the same notation as in the proof of Theorem \ref{thm-NS}. For any fixed $\delta>0$ and $\alpha\in \T$, set 
$$X_n:= \bigg(\frac{1-|\lambda|^2}{2}\bigg)^{-1/2} \cdot {\rm Re}\, \bigg( \frac{\alpha b_{N+1-n} Y_{N+1-n}}{\rho_N}\bigg), \quad 1\leq n \leq N.$$ Then we have
\begin{align}\label{eq-2+2delta-norm}
	\sum_{n=1}^N \E\big[ |X_n|^{2+2\delta}\big]\lesssim_{\delta,\lambda} \frac{\sum_{n=1}^N |b_n|^{2+2\delta}}{\rho_N^{2+2\delta}}
	\lesssim_{\delta,\lambda} \bigg( \sum_{n=1}^N |a_n|^{2+2\delta}\bigg) \bigg( \sum_{n=1}^N |a_n|^2\bigg)^{-1-\delta}.
\end{align}
Moreover, dividing both sides of \eqref{eq-VN-value} by $(1-|\lambda|^2)/2$, we obtain
\begin{align*} 
	\sum_{n=1}^N \E\big[ X_n^2 | \cal F_{n-1}\big]= 1+ \bigg(\frac{1-|\lambda|^2}{2}\bigg)^{-1} \frac{1}{\rho^2_N} \sum_{n=1}^N \bigg(\frac{{\rm Re}\,(\alpha^2 \mu \,b_n^2f^{\circ (n+1)})}{2}\bigg).
\end{align*}
This implies 
\begin{align*}
	&\bigg|\sum_{n=1}^N \E\big[ X_n^2 | \cal F_{n-1}\big] -1\bigg|^{1+\delta}\lesssim_{\delta,\lambda} \bigg| \frac{1}{\rho_N^2}\sum_{n=1}^N b_n^2 f^{\circ (n+1)}\bigg|^{1+\delta}.
\end{align*}
Now taking the expectation and using a Khintchine type inequality \cite[Corollary 1.4]{Nic22}, we get
\begin{align*} 
	\E\bigg[\, \bigg|\sum_{n=1}^N \E\big[ X_n^2 | \cal F_{n-1}\big] -1\bigg|^{1+\delta}\bigg] \lesssim_{\delta,f}& \bigg(\sum_{n=1}^N |b_n|^4 \bigg)^{\frac{1+\delta}{2}} \rho_N^{-2-2\delta} \\ 
	\lesssim_{\delta,f} &\bigg( \sum_{n=1}^N |a_n|^4\bigg)^{\frac{1+\delta}{2}}\bigg( \sum_{n=1}^N |a_n|^2\bigg)^{-1-\delta}
\end{align*}

Combining this with \eqref{eq-2+2delta-norm}, we deduce from Haeusler's theorem that
\begin{align}\label{eq-using-Haeusler}
\sup_{x\in \R}\bigg| \PP\Big( \sum_{n=1}^N X_n\leq x \Big) - \PP \Big(\cal N(0,1) \leq x \Big)\bigg| 
\lesssim_{\delta,f} \bigg(\sum_{n=1}^N |a_n|^{4} \bigg)^{\frac{1+\delta}{6+4\delta}}  \bigg(\sum_{n=1}^N |a_n|^{2} \bigg)^{-\frac{1+\delta}{3+2\delta}}
\end{align} whenever $\delta\geq 1$.


Moreover, by the identity \eqref{eq-sum-to-martingale}, standard calculations show that
\begin{align*} 
	\sum_{n=1}^N X_n = \sqrt{2} \cdot \bigg( p_N {\rm Re}\,\big( \alpha Z_N \big)- R_N\bigg),
\end{align*}
where $$p_N: =\frac{\sigma_N}{{\sqrt{1-|\lambda|^2}\rho_N}}$$ and the remainder $R_N$ satisfies $$|R_N| \leq \frac{|\lambda|}{\sqrt{1-|\lambda|^2}} \frac{|b_N|}{\rho_N} =: q_N.$$ Consequently, the left-hand side of \eqref{eq-using-Haeusler} can be rewritten as
$$\sup_{x\in \R} \bigg| \PP\bigg( {\rm Re}\, \big( \alpha Z_N\big)\leq \frac{x-R_N}{p_N} \bigg) - \PP\bigg( \cal N(0,1/2) \leq x \bigg)\bigg|.$$

Since $\lim_{N\to \infty}p_N=1$, we see for sufficiently large $N$, 
$$\PP\bigg( {\rm Re}\, \big( \alpha Z_N\big)\leq \frac{x-q_N}{p_N} \bigg) \leq \PP\bigg( {\rm Re}\, \big( \alpha Z_N\big)\leq \frac{x-R_N}{p_N} \bigg) \leq \PP\bigg( {\rm Re}\, \big( \alpha Z_N\big)\leq \frac{x+q_N}{p_N} \bigg).$$ 

Now applying Lemma \ref{lem-normal-ball-probability}, we conclude that 
\begin{align*} 
	& \sup_{x\in \R} \bigg[ \PP\bigg( {\rm Re}\, \big( \alpha Z_N\big)\leq x \bigg)- \PP\bigg( \cal N(0,1/2) \leq x \bigg) \bigg] \\
	= & \sup_{x\in \R} \bigg[ \PP\bigg( {\rm Re}\, \big( \alpha Z_N\big)\leq \frac{x-q_N}{p_N} \bigg)- \PP\bigg( \cal N(0,1/2) \leq \frac{x-q_N}{p_N} \bigg) \bigg] \\
	\leq &  \sup_{x\in \R} \bigg[ \PP\bigg( {\rm Re}\, \big( \alpha Z_N\big)\leq \frac{x-R_N}{p_N} \bigg)- \PP\bigg( \cal N(0,1/2) \leq \frac{x-q_N}{p_N} \bigg) \bigg]\\
	\leq  & \sup_{x\in \R} \bigg[ \PP\bigg( {\rm Re}\, \big( \alpha Z_N\big)\leq \frac{x-R_N}{p_N} \bigg)- \PP\bigg( \cal N(0,1/2) \leq x \bigg) \bigg] \\
	& + \sup_{x\in \R} \bigg[ \PP\bigg( \cal N(0,1/2) \leq x \bigg)- \PP\bigg( \cal N(0,1/2) \leq \frac{x-q_N}{p_N} \bigg) \bigg].
\end{align*}
The first supremum can be bounded using \eqref{eq-using-Haeusler}. Next we consider the second one. By \eqref{eq-rhoN-sigmaN} we have $\sigma_N^2 = (1-|\lambda|^2) \rho^2_N +|\lambda|^2 |b_N|^2,$ and hence $$|p_N-1|\leq \frac{|\lambda|}{\sqrt{1-|\lambda|^2}} \frac{|b_N|}{\rho_N}.$$ Now an application of Lemma \ref{lem-normal-ball-probability} gives 
\begin{align*} 
	& \sup_{x\in \R} \bigg[ \PP\bigg( \cal N(0,1/2) \leq x \bigg)- \PP\bigg( \cal N(0,1/2) \leq \frac{x-q_N}{p_N} \bigg) \bigg] \\
	 =& \sup_{x\in \R} \bigg[ \PP\bigg( \cal N(0,1/2) \leq p_Nx+q_N \bigg)- \PP\bigg( \cal N(0,1/2) \leq x \bigg) \bigg]\\
	 \leq & 1.35\cdot (|p_N-1|+|q_N|) \\
	 \leq & 2.7\cdot \frac{|\lambda|}{\sqrt{1-|\lambda|^2}} \frac{|b_N|}{\rho_N}.
\end{align*}
Then we can apply \eqref{eq-2-norm-eqvalence} and \eqref{eq-rhoN-sigmaN} to obtain
\[
 \sup_{x\in \R} \bigg[ \PP\bigg( \cal N(0,1/2) \leq x \bigg)- \PP\bigg( \cal N(0,1/2) \leq \frac{x-q_N}{p_N} \bigg) \bigg]\leq 2.7\cdot \frac{|\lambda|}{1-|\lambda|} \frac{|b_N|}{\sqrt{\sum_{n=1}^N |a_n|^2}}.
\]
Combining these estimates yields
\begin{align*} 
	&\sup_{x\in \R} \bigg[ \PP\Big( {\rm Re}\, \big( \alpha Z_N\big) \leq x \Big) - \PP\Big( \cal{N}(0,1/2) \leq x \Big)\bigg] \\
	\leq & \,\,C(\delta,f) \cdot \bigg(\sum_{n=1}^N |a_n|^{4} \bigg)^{\frac{1+\delta}{6+4\delta}}  \bigg(\sum_{n=1}^N |a_n|^{2} \bigg)^{-\frac{1+\delta}{3+2\delta}} + 2.7 \cdot \frac{|\lambda|}{1-|\lambda|} \frac{\big|\sum_{n=1}^N \lambda^{N-n} a_n\big|}{\sqrt{\sum_{n=1}^N |a_n|^2}}.
\end{align*}
The same upper bound for $$\sup_{x\in \R} \bigg[ \PP\Big( \cal{N}(0,1/2) \leq x \Big)-\PP\Big( {\rm Re}\, \big( \alpha Z_N\big) \leq x \Big) \bigg]$$ can be obtained by a similar argument, and the proof of Theorem \ref{thm-BE-inner} is completed.
\end{proof}



\section{Proofs of Lemma \ref{lem-Yn} and Lemma \ref{lem-transfer}}



	\begin{proof}[Proof of Lemma \ref{lem-Yn}] Clearly, $Y_n=f^{\circ n} - \lambda f^{\circ (n+1)}$ is $\F_{n}$-measurable. Furthermore, since $f^{\circ n}$ preserves the Lebesgue measure, one has
			\begin{align}
				\big\la f^{\circ n}, \big(f^{\circ (n+1)}\big)^k \big\ra = \big\la z \circ f^{\circ n}, f^k \circ f^{\circ n}\big\ra= \big\la z, f^k\big\ra =
				\begin{cases}
					\lambda\,, &  \text{if }k=1; \\ 
					 0\,, &  \text{otherwise}.
				\end{cases}
		\end{align}
		Thus $$\E\big[ f^{\circ n} | \F_{n+1}\big] = \sum_{k\in \mathbb Z} \big\la f^{\circ n}, \big(f^{\circ (n+1)}\big)^k \big\ra \cdot \big(f^{\circ (n+1)}\big)^k = \lambda f^{\circ (n+1)}.$$ This shows $(Y_n)_{n=1}^\infty$ is indeed a reverse martingale difference sequence. 

		Since $\overline{f^{\circ (n+1)}}$ is $\mathscr{F}_{n+1}$-measurable, we have $$\E\big[ f^{\circ n} \ol{f^{\circ (n+1)}} | \F_{n+1}\big] = \ol{f^{\circ (n+1)}} \cdot \E\big[ f^{\circ n} | \F_{n+1}\big] = \ol{f^{\circ (n+1)}}\cdot \lambda f^{\circ (n+1)} = \lambda.$$ This implies that 
		\begin{align*} 
			\E\big[|Y_n|^2 | \F_{n+1}\big] =& \E\big[ 1+ |\lambda|^2 - \lambda \ol{f^{\circ n}} f^{\circ (n+1)} - \ol{\lambda} f^{\circ n} \ol{f^{\circ (n+1)}} | \F_{n+1}\big]\\
			=&1+ |\lambda|^2  - \lambda \E\big[ \,\ol{f^{\circ n}} f^{\circ (n+1)}|\F_{n+1}\big] - \ol{\lambda} \E\big[f^{\circ n} \ol{f^{\circ (n+1)}} | \F_{n+1}\big]\\
			=& 1-|\lambda|^2.
		\end{align*}
		The identity in \ref{item:|Yn|2} follows.

		By expanding the square, we obtain $$Y_n^2 = \big(f^{\circ n}\big)^2- 2 \lambda f^{\circ n} f^{\circ (n+1)} + \lambda^2 \big(f^{\circ (n+1)}\big)^2.$$ 
		Note that the conditional expectation of the cross term is $$\E\big[2 \lambda f^{\circ n} f^{\circ (n+1)} | \F_{n+1}\big] =  2 \lambda^2  \big(f^{\circ (n+1)}\big)^2 .$$
		Moreover, it holds that 
		\begin{align*} 
			\big\la \big(f^{\circ n}\big)^2 , \big(f^{\circ (n+1)}\big)^k \big\ra = \big\la z^2 \circ f^{\circ n}, f^k \circ f^{\circ n}\big\ra= \big\la z^2 , f^k\big\ra =
				\begin{cases}
					\mu\, , &   \text{if }k=1; \\ 
					 \lambda^2 \, , &  \text{if }k=2;\\
					 0\, , & \text{otherwise}.
				\end{cases}
		\end{align*}
		This implies 
		\begin{align*} 
			\E\big[ \big(f^{\circ n}\big)^2  | \F_{n+1}\big] =& \sum_{k\in \mathbb Z} \big\la \big( f^{\circ n}\big)^2 , \big(f^{\circ (n+1)}\big)^k \big\ra \cdot \big(f^{\circ (n+1)}\big)^k\\ 
			=&\mu f^{\circ (n+1)} + \lambda^2 \big(f^{\circ (n+1)}\big)^2.
	\end{align*}
	Now \ref{item:Yn2} follows by combining these identities.
	
	The identity in \ref{item:reYn2} is a direct consequence of \ref{item:|Yn|2} and \ref{item:Yn2}. Indeed, observe that $$|\alpha Y_n|^2 = \big( {\rm Re}\, \alpha Y_n \big)^2 + \big({\rm Im}\, \alpha Y_n \big)^2, \quad \text{and}\quad {\rm Re}\, (\alpha Y_n)^2 = \big( {\rm Re}\, \alpha Y_n \big)^2 - \big({\rm Im}\, \alpha Y_n \big)^2.$$
	By taking the conditional expectation, we get 
	\begin{align*} 
		2 \E\big[({\rm Re}\, \alpha Y_n)^2 |\mathscr F_{n+1}\big]=& \E\big[ |\alpha Y_n|^2 | \F_{n+1}\big]+ \E\big[  {\rm Re}\, \big(\alpha Y_n \big)^2 | \F_{n+1}\big]\\
		=& |\alpha |^2 \,\E\big[ |Y_n|^2 | \F_{n+1}\big]+ {\rm Re}\, \big( \alpha^2 \, \E\big[ Y_n^2 | \F_{n+1}\big]\big)\\
		=& |\alpha|^2\,(1-|\lambda|^2)+ {\rm Re}\,(\alpha^2 \mu f^{\circ (n+1)}).
	\end{align*}
	The proof of Lemma \ref{lem-Yn} is completed.
	\end{proof}


		\begin{proof}[Proof of Lemma \ref{lem-transfer}]

		The claim \ref{item:keep-no-ell2} follows trivially from
		\begin{align*}
			a_1=b_1,\,\, \text{and} \,\, a_n=b_n-\lambda b_{n-1}, \text{ for }n\geq 2.
		\end{align*}
		Moreover, the claim \ref{item:keep-ell2} can be deduced by taking the $L^2(\T)$ norm on both sides of the following identity:
		$$\sum_{n=1}^\infty a_n z^n = (1-\lambda z) \cdot \bigg( \sum_{n=1}^\infty b_nz^n \bigg).$$
		Next we prove the claim \ref{item:keep-Lindberg}. Consider the generating functions $$A_N(z):=\sum_{n=1}^N a_n z^n,\quad \text{and} \quad B_N(z):=\sum_{n=1}^N b_n z^n.$$
		Then for $N\geq 2$ we have 
		\begin{align*}
			A_N(z)=(1-\lambda z) B_N(z) + \lambda b_Nz^{N+1}.
		\end{align*}
		This implies 
		\begin{align*}
		\big\|A_N \big\|_{L^2(\T)} \leq  \big(1+|\lambda|\big)\big\|B_N \big\|_{L^2(\T)}+|\lambda| |b_N|.
	\end{align*}
	
   Substituting $$|b_N| \leq \sum_{n=1}^N |\lambda|^{N-n} |a_n| \leq  \frac{\max_{1\leq n \leq N} |a_n|}{1-|\lambda|},$$ and using Parseval's identity, one has
	\begin{align}\label{eq1.2}
		\sqrt{\sum_{n=1}^N |a_n|^2} \leq (1+|\lambda|) \sqrt{\sum_{n=1}^N |b_n|^2} + \frac{ \lambda }{1-|\lambda|}\max_{1\leq n \leq N} |a_n|.
	\end{align}
	
	Now if $$\lim_{N\to \infty} \frac{\max_{1\leq n \leq N} |a_n|}{\sqrt{\sum_{n=1}^N |a_n|^2}} = 0,$$ then from \eqref{eq1.2} we deduce that 
	\begin{align}\label{eq-ratio-lower-bound}
		\liminf_{N\to \infty} \frac{\sqrt{\sum_{n=1}^N |b_n|^2}}{\sqrt{\sum_{n=1}^N |a_n|^2}}\geq \frac{1}{1+|\lambda|}.
	\end{align}
	This implies $$\limsup_{N\to \infty} \frac{|b_N|}{\sqrt{\sum_{n=1}^N |b_n|^2}}\leq \limsup_{N\to \infty}\frac{1+|\lambda|}{1-|\lambda|} \cdot \frac{\max_{1\leq n \leq N} |a_n|}{\sqrt{\sum_{n=1}^N |a_n|^2}} =0.$$ The claim \ref{item:keep-Lindberg} follows.
	\end{proof} 








	
	%----------------------------------Text of article-------------------------------------------
	
	
	
	
	
	
	
	
	
	
	
	
	
	
	
	
	
	
	
	
	
	
	%-------------------------------------reference-----------------------------------------



\begin{thebibliography}{FGL00}

		\bibitem{Bol82}
		E. Bolthausen, \emph{Exact convergence rates in some martingale central limit theorems.} {\sl Ann. Prob.} {\bf 10} (1982), 672--688.
		
		\bibitem{BE71}
		 B. M. Brown and G. K. Eagleson, \emph{Martingale convergence to infinitely divisible laws with finite variances.} {\sl Trans. Amer. Math. Soc.} {\bf 162} (1971), 449--453. 


		\bibitem{DM91}
		C.I. Doering, R. Ma\~n\'e, \emph{The dynamics of inner functions.} {\sl Ens. Mat.} {\bf 3} (1991) 5--79.

		
		 \bibitem{Eag75}
		 G. K. Eagleson, \emph{Martingale convergence to mixtures of infinitely divisible laws.} {\sl Ann. Prob.} {\bf 3} (1975), 557--562. 

		\bibitem{Gar81}
		J. B. Garnett, {\sl Bounded Analytic Functions.} Academic Press, New York, 1981.

		\bibitem{Hae88}
		E. Haeusler, \emph{On the rate of convergence in the central limit theorem for martingales with discrete and continuous time.} {\sl Ann. Prob.} {\bf 16}, 275--299.

		\bibitem{Hof62}
		K. Hoffman, {\sl Banach Spaces of Analytic Functions.} Prentice-Hall, Englewood Cliffs, NJ, 1962.

		\bibitem{IU25}
		O. Ivrii, M. Urba\'nski, \emph{Inner Functions, Composition Operators, Symbolic Dynamics and Thermodynamic Formalism.} preprint, arXiv:2308.16063v2.

		\bibitem{Loy69}
		  R. M. Loynes, \emph{The central limit theorem for backwards martingales.} {\sl Z. Wahrscheinlichkeitstheorie und Verw. Gebiete}, {\bf 13} (1969), 1--8.

		\bibitem{Nic22}
		A. Nicolau, \emph{Convergence of linear combinations of iterates of an inner function.} {\sl J. Math. Pures Appl.}, {\bf 161} (2022), 135--165.
		

		\bibitem{NS22}
		A. Nicolau and O. Soler i Gibert, \emph{A central limit theorem for inner functions.} {\sl Adv. Math.} {\bf 401} (2022), Paper No. 108318, 39 pages.

		\bibitem{NS24}
		A. Nicolau and O. Soler i Gibert, \emph{The Central Limit Theorem for inner functions II}. (2026). In: {\sl New Trends in Complex Analysis, Fourier Analysis, and Operator Theory.} Springer INdAM Series, vol 66. Springer, Singapore.

		\bibitem{Pom81}
		C. Pommerenke, \emph{On ergodic properties of inner functions.} {\sl Math. Ann.} {\bf 256} (1981) 43--50.

		
		
		

		
		
	\end{thebibliography}
	
	
	
	
	
	
	
	
\end{document}
